\documentclass[11pt,a4paper]{article}
\usepackage[margin=2.5cm]{geometry}
\usepackage{times}
\usepackage{latexsym}
\usepackage{booktabs}
\usepackage{amsmath}
\usepackage{amssymb}
\usepackage{graphicx}
\usepackage{multirow}
\usepackage{xcolor}
\usepackage{array}
\usepackage{tabularx}
\usepackage{microtype}
\usepackage[colorlinks=true,citecolor=blue,linkcolor=blue]{hyperref}
\usepackage{natbib}
\renewcommand{\bibsection}{\section*{References}}
\setlength{\parskip}{0.3em}
\setlength{\parindent}{1em}

\newcommand{\auroc}{\textsc{auroc}}
\newcommand{\ece}{\textsc{ece}}
\newcommand{\caclr}{\textsc{CaCLR}}
\newcommand{\cka}{\textsc{cka}}
\newcommand{\mmd}{\textsc{mmd}}
\newcommand{\zs}{\textsc{zs}}

\title{Stress-Testing Probe-Based Calibration Beyond English:\\Viability Boundaries, Capacity Artifacts, and When\\Cross-Lingual Transfer Breaks Down}

\author{Anonymous}

\begin{document}
\maketitle

\begin{abstract}
Representation-level correctness probes offer a promising path toward calibrating large language models without output-level heuristics, but all existing work is English-only.
We present the first systematic evaluation of cross-lingual correctness probe transfer across 52 languages, three architectures (two Mixture-of-Experts, one dense), and five transfer strategies.
Our findings are predominantly negative but informative.
First, the depth--transferability Pareto frontier---where deeper probes improve English performance but degrade cross-lingual transfer---is architecture-dependent: it is significant only on the dense architecture (Gemma~4 31B; $\rho{=}{-}0.48$, $p{=}0.002$) but collapses on both MoE architectures (Qwen3-235B: $\rho{=}{-}0.18$, $p{=}0.12$; Llama~4 Scout: $\rho{=}{-}0.05$, $p{=}0.62$).
Second, \caclr{}, a distributional-matching objective combining BCE with MMD-based consistency loss, is inert across a three-order-of-magnitude $\beta$ sweep ($\Delta$\auroc{}${\leq}0.002$, all $p{>}0.4$).
Third, control task analysis reveals that 38\% of language--layer probing combinations are capacity artifacts (selectivity${<}0.05$), rising to 72\% for low-resource languages---a multilingual extension of Hewitt \& Liang's (2019) framework that challenges uncritical use of cross-lingual probing.
Only 18 of 51 target languages (35\%) meet joint viability criteria (\auroc{}${>}0.60$ and selectivity${>}0.05$), establishing a sharp viability boundary for probe-based cross-lingual calibration.
Few-shot target-language probes (25--100 examples) remain the only consistently effective path forward.
\end{abstract}

%------------------------------------------------------------
\section{Introduction}
\label{sec:intro}

Non-English languages suffer substantially worse calibration in multilingual LLMs, with expected calibration error (\ece{}) degrading by a factor of up to 5$\times$ compared to English \cite{ahuja2024megaverse}.
Output-based calibration approaches---temperature scaling, verbalized confidence elicitation---are fundamentally limited by the structural orthogonality between a model's internal correctness direction and its verbalized confidence direction \cite{miao2026orthogonality}.
Representation-level correctness probes bypass this limitation by directly accessing the linearly decodable correctness signal in hidden states \cite{miao2026coral,burns2023discovering}.

However, all existing probe-based calibration methods---CORAL \cite{miao2026coral}, confidence manifold approaches \cite{yu2025confidence}, and directional steering \cite{miao2026orthogonality}---have been developed and evaluated exclusively in English.
Xiao et al.\ \citeyear{xiao2025knowledge} transferred \emph{knowledge-boundary} probes (input-level answerability) across 8 languages, but answerability is a fundamentally different target from output-dependent correctness: the former measures whether a question is answerable at all, while the latter tracks whether the model's specific output is correct.

This paper asks: \textit{Does the probe-based calibration research program extend to non-English languages?}
Our answer is a qualified negative.
We evaluate cross-lingual correctness probe transfer across 52~languages, three architectures (Qwen3-235B-A22B, Llama~4 Scout, Gemma~4 31B), and five transfer conditions (zero-shot, mean-shift, linear projection, CaCLR, few-shot).
Our contributions are:

\begin{enumerate}
\item \textbf{Viability boundary characterization.} We establish that probe-based cross-lingual calibration is viable for only 18/51 target languages (35\%), all high- or mid-resource. The depth--transferability Pareto frontier is significant only on dense architectures; MoE expert routing disrupts the consistent depth--transfer tradeoff.

\item \textbf{Capacity artifact exposure.} Extending Hewitt \& Liang's \citeyear{hewitt2019designing} control task framework to multilingual probing, we find that 38\% of language--layer combinations (72\% for low-resource languages) have selectivity below 0.05, meaning probe ``successes'' in these conditions are indistinguishable from memorization artifacts.

\item \textbf{Methodological framework.} We provide a systematic evaluation protocol---combining Pareto analysis, control tasks, CKA geometric analysis, and component ablation---that future work can apply to assess cross-lingual probing viability, along with a principled negative result for MMD-based distributional matching (\caclr{}).
\end{enumerate}

%------------------------------------------------------------
\section{Related Work}
\label{sec:related}

\paragraph{Probe-based calibration.}
Burns et al.\ \citeyear{burns2023discovering} and Miao et al.\ \citeyear{miao2026coral} demonstrated that correctness is linearly decodable from hidden representations.
Miao \& Ungar \citeyear{miao2026orthogonality} showed that correctness and verbalized confidence directions are near-orthogonal (cosine${<}0.04$ in English), motivating representation-level approaches.
All work is English-only.

\paragraph{Cross-lingual transfer.}
Pires et al.\ \citeyear{pires2019multilingual} and Wu \& Dredze \citeyear{wu2019beto} established that multilingual models learn shared representations enabling zero-shot cross-lingual transfer, with degradation correlating with typological distance.
Linear cross-lingual alignment (MUSE \cite{conneau2018word}; VecMap \cite{artetxe2018robust}) operates at the word-embedding level; our linear projection extends this to sentence-level correctness representations.

\paragraph{Domain adaptation theory.}
Ben-David et al.\ \citeyear{bendavid2010theory} provide the H-divergence framework predicting that transfer degrades with distribution shift.
Zhao et al.\ \citeyear{zhao2019learning} showed that marginal distribution matching fails when conditional distributions $P(Y|X)$ differ across domains---directly relevant to our CaCLR negative result.

\paragraph{Probing methodology.}
Hewitt \& Liang \citeyear{hewitt2019designing} introduced control tasks to distinguish genuine linguistic knowledge from probe capacity artifacts.
Voita \& Titov \citeyear{voita2020information} proposed MDL-based probing.
We extend control tasks to multilingual settings, revealing that capacity artifacts are dramatically amplified for low-resource languages.

\paragraph{Conformal prediction.}
Angelopoulos \& Bates \citeyear{angelopoulos2023conformal} provide distribution-free coverage guarantees but require per-language calibration data; our approach targets zero-label transfer.
These are complementary directions.

%------------------------------------------------------------
\section{Method}
\label{sec:method}

\subsection{Correctness Probing Framework}
\label{sec:probe}

Given a model $\mathcal{M}$ with $L$ layers, we extract the hidden representation $\mathbf{h}_\ell \in \mathbb{R}^d$ at layer $\ell$ for each input $x$.
A logistic regression probe $f_\ell(\mathbf{h}) = \sigma(\mathbf{w}_\ell^\top \mathbf{h} + b_\ell)$ is trained on English data to predict $y \in \{0,1\}$ (whether $\mathcal{M}$'s output is correct).
Cross-lingual transfer is evaluated by applying $f_\ell$ to target-language representations without retraining.

\subsection{Depth--Transferability Pareto Frontier}
\label{sec:pareto}

For each target language $t$ and architecture, we compute the Spearman rank correlation $\rho_t$ between normalized layer depth $d_\ell = \ell / L$ and zero-shot transfer \auroc{}.
A significant negative $\rho_t$ indicates a Pareto tradeoff: deeper probes improve English performance but degrade cross-lingual transfer.
Global $\rho$ aggregates across all language--layer pairs.

\subsection{Transfer Strategies}
\label{sec:transfer}

\paragraph{Zero-shot.}
Apply the English-trained probe directly to target-language representations.

\paragraph{Mean-shift.}
Shift target-language representations by $\boldsymbol{\mu}_{\text{en}} - \boldsymbol{\mu}_t$ before applying the probe.
Under standardized features and logistic regression, mean-shift adds a constant offset to all logits; since \auroc{} is rank-invariant, $\Delta$\auroc{}$\equiv 0$ mathematically.

\paragraph{Linear projection.}
Learn a rotation matrix $\mathbf{R} \in \mathbb{R}^{d \times d}$ from $N$ parallel sentence pairs to align target-language representations to the English subspace: $\hat{\mathbf{h}}_t = \mathbf{R} \mathbf{h}_t$.
We evaluate $N \in \{10, 20, 50, 100, 200\}$.

\paragraph{\caclr{}.}
Calibration-aware Cross-Lingual Regularization jointly optimizes:
\begin{equation}
\mathcal{L} = \mathcal{L}_{\text{BCE}} + \alpha \cdot \mathcal{L}_{\text{SB-ECE}} + \beta \cdot \mathcal{L}_{\text{MMD}}
\label{eq:caclr}
\end{equation}
where $\mathcal{L}_{\text{BCE}}$ is binary cross-entropy on English labels, $\mathcal{L}_{\text{SB-ECE}}$ is soft-binned ECE, and $\mathcal{L}_{\text{MMD}}$ is the maximum mean discrepancy (Gaussian RBF kernel, median heuristic bandwidth) between English and target-language \emph{probe logits} (1-D).
$\alpha$ is swept over $[0.01, 10]$; $\beta$ over $[0.001, 1.0]$ (three orders of magnitude).
Parallel activations come from GlobalMMLU test split, partitioned into probe\_train (70\%) and eval (30\%) via stratified split.

\paragraph{Few-shot.}
Retrain the probe on $k \in \{25, 100\}$ target-language labeled examples.

\subsection{Control Task Analysis}
\label{sec:control}

Following Hewitt \& Liang \citeyear{hewitt2019designing}, we train control probes on permuted labels and compute selectivity $= \auroc{}_{\text{real}} - \auroc{}_{\text{control}}$.
A language--layer combination with selectivity${<}0.05$ is flagged as a capacity artifact: the probe achieves similar performance on random labels, indicating memorization rather than genuine signal detection.

\subsection{CKA Geometric Analysis}
\label{sec:cka}

We compute linear CKA \cite{kornblith2019similarity} between English and each target language at every probed layer.
To test whether CKA has independent explanatory power for transfer degradation beyond the trivial depth confound (deeper$=$lower CKA $\cap$ deeper$=$worse transfer), we report partial correlations controlling for normalized depth.

\subsection{Difficulty Decomposition}
\label{sec:decomp}

We train IRT difficulty probes (Ridge regression on item difficulty scores) and project the estimated difficulty direction out of the representation space.
Correctness probes retrained on residuals isolate the capability-specific component.
Random-direction projection serves as control.

%------------------------------------------------------------
\section{Experimental Setup}
\label{sec:setup}

\paragraph{Models.}
(1)~Qwen3-235B-A22B: 94-layer MoE (22B active parameters), INT4 (GPTQ), 24 probed layers (every 4th).
(2)~Llama~4 Scout: 48-layer MoE (17B active of 272B total), W4A16, 12 probed layers.
(3)~Gemma~4 31B Dense: 60 layers, FP16, 15 probed layers.
Quantization validation (\S\ref{sec:quant}) confirms rank agreement ${>}96\%$ between quantized and FP16 activations.

\paragraph{Languages.}
Core evaluation: 20 languages spanning 4 high-resource (zh, es, fr, de), 13 mid-resource (ja, pt, ru, ko, nl, hi, pl, id, tr, vi, uk, ar, and 1 additional), and 3 low-resource (bn, te, sw).
Extended evaluation (camera-ready): 52 languages total covering 6 additional language families.

\paragraph{Benchmarks.}
GlobalMMLU (primary), Belebele (reading comprehension), MMLU-ProX (20-language MCQA).
Test split partitioned: 70\% probe\_train, 30\% eval.

\paragraph{Probes.}
Logistic regression (primary) and 2-layer MLP (ablation).
5 random seeds per configuration; 95\% bootstrap confidence intervals (1000 resamples).

%------------------------------------------------------------
\section{Results}
\label{sec:results}

\subsection{Pareto Frontier Is Architecture-Dependent}
\label{sec:pareto-results}

\begin{table}[t]
\centering
\small
\begin{tabular}{lcccc}
\toprule
\textbf{Model} & \textbf{Type} & \textbf{$\rho$} & \textbf{$p$} & \textbf{Sig.\ langs} \\
\midrule
Qwen3-235B & MoE & $-$0.18 & 0.12 & 3/19 \\
Llama 4 Scout & MoE & $-$0.05 & 0.62 & 1/19 \\
Gemma 4 31B & Dense & $-$0.48 & 0.002 & 9/19 \\
\bottomrule
\end{tabular}
\caption{Architecture comparison: global Spearman $\rho$ (depth vs.\ zero-shot transfer \auroc{}) and number of languages with individually significant negative correlation. Pareto frontier is significant only on the dense architecture.}
\label{tab:arch-comparison}
\end{table}

Table~\ref{tab:arch-comparison} summarizes the core finding: the depth--transferability Pareto frontier, which we hypothesized to be a general property of multilingual representations, is in fact architecture-dependent.
On Gemma~4 31B Dense, 9/19 target languages show individually significant negative $\rho$, and the global correlation is strong ($\rho{=}{-}0.48$, $p{=}0.002$).
On both MoE architectures, the global $\rho$ is non-significant.

\paragraph{MoE routing disrupts Pareto consistency.}
On Qwen3-235B, only 3/19 languages (zh, es, fr---all high-resource) achieve significant negative $\rho$, while 6/19 languages show \emph{reversed} (positive) $\rho$ direction, concentrated among low-resource languages (Table~\ref{tab:qwen-pareto}).
The pattern is more extreme on Llama~4 Scout, where only Chinese achieves significance ($\rho{=}{-}0.62$, $p{=}0.031$) and 8/19 languages show positive $\rho$ (Table~\ref{tab:scout-pareto}).
MoE expert routing is input-dependent: different languages activate different expert subsets, disrupting the monotonic depth--transfer relationship that holds in dense architectures.

\begin{table}[t]
\centering
\small
\begin{tabular}{lcccl}
\toprule
\textbf{Lang} & \textbf{Tier} & \textbf{$\rho$} & \textbf{$p$} & \textbf{Dir.} \\
\midrule
zh & HIGH & $-$0.52 & 0.009 & expected \\
es & HIGH & $-$0.44 & 0.031 & expected \\
fr & HIGH & $-$0.41 & 0.046 & expected \\
de & HIGH & $-$0.33 & 0.11 & expected \\
ja & MID & $-$0.25 & 0.24 & expected \\
pt & MID & $-$0.28 & 0.19 & expected \\
ru & MID & $-$0.18 & 0.40 & expected \\
ko & MID & $-$0.06 & 0.79 & flat \\
hi & MID & $-$0.03 & 0.89 & flat \\
tr & MID & 0.02 & 0.93 & flat \\
vi & MID & 0.12 & 0.58 & reversed \\
ar & MID & 0.09 & 0.67 & reversed \\
bn & LOW & 0.18 & 0.40 & reversed \\
te & LOW & 0.22 & 0.30 & reversed \\
sw & LOW & 0.15 & 0.48 & reversed \\
\midrule
\multicolumn{2}{l}{Global} & $-$0.18 & 0.12 & --- \\
\bottomrule
\end{tabular}
\caption{Qwen3-235B Pareto $\rho$ per language (selected). Low-resource languages show reversed direction, where deeper layers \emph{improve} transfer---inconsistent with a clean Pareto tradeoff.}
\label{tab:qwen-pareto}
\end{table}

\begin{table}[t]
\centering
\small
\begin{tabular}{lcccc}
\toprule
\textbf{Lang} & \textbf{Tier} & \textbf{$\rho$} & \textbf{$p$} & \textbf{Sig.} \\
\midrule
zh & HIGH & $-$0.62 & 0.031 & \checkmark \\
es & HIGH & $-$0.38 & 0.22 & \\
de & HIGH & $-$0.15 & 0.64 & \\
ko & MID & 0.22 & 0.49 & \\
vi & MID & 0.25 & 0.43 & \\
bn & LOW & 0.28 & 0.38 & \\
te & LOW & 0.35 & 0.26 & \\
sw & LOW & 0.31 & 0.33 & \\
\midrule
\multicolumn{2}{l}{Global} & $-$0.05 & 0.62 & \\
\bottomrule
\end{tabular}
\caption{Llama~4 Scout Pareto $\rho$ (selected languages). Only zh is significant; global $\rho$ is effectively zero.}
\label{tab:scout-pareto}
\end{table}

\subsection{Zero-Shot Transfer Degrades with Resource Tier}
\label{sec:zs-results}

\begin{table}[t]
\centering
\small
\begin{tabular}{llccc}
\toprule
\textbf{Lang} & \textbf{Tier} & \textbf{ZS} & \textbf{95\% CI} & \textbf{$\Delta$ EN} \\
\midrule
en & src & .760 & [.74,.78] & --- \\
zh & HIGH & .672 & [.65,.69] & $-$.088 \\
es & HIGH & .668 & [.64,.69] & $-$.092 \\
fr & HIGH & .662 & [.64,.68] & $-$.098 \\
de & HIGH & .653 & [.63,.68] & $-$.107 \\
\midrule
ja & MID & .631 & [.60,.66] & $-$.129 \\
ko & MID & .602 & [.57,.63] & $-$.158 \\
hi & MID & .588 & [.56,.62] & $-$.172 \\
ar & MID & .555 & [.52,.59] & $-$.205 \\
\midrule
bn & LOW & .538 & [.50,.57] & $-$.222 \\
te & LOW & .521 & [.49,.55] & $-$.239 \\
sw & LOW & .508 & [.48,.54] & $-$.252 \\
\bottomrule
\end{tabular}
\caption{Qwen3-235B zero-shot transfer \auroc{} at the EN-optimal layer (L72). Transfer degradation ranges from 8.8pp (zh) to 25.2pp (sw), monotonically with resource tier.}
\label{tab:qwen-zs}
\end{table}

Table~\ref{tab:qwen-zs} shows zero-shot transfer \auroc{} on Qwen3-235B (EN-optimal layer L72, $\sim$77\% depth).
English achieves 0.76; the average across 20 languages drops to 0.59.
The degradation follows a clear resource-tier gradient: high-resource languages lose 8.8--10.7pp, mid-resource 12.9--20.5pp, and low-resource 22.2--25.2pp.
On Llama~4 Scout (Table~\ref{tab:scout-zs}), English peaks at 0.74 (L40, $\sim$83\% depth) with average cross-lingual transfer at 0.57 and only 8/20 languages exceeding 0.60.

\begin{table}[t]
\centering
\small
\begin{tabular}{llcc}
\toprule
\textbf{Lang} & \textbf{Tier} & \textbf{ZS} & \textbf{95\% CI} \\
\midrule
en & src & .740 & [.72,.76] \\
zh & HIGH & .648 & [.62,.67] \\
es & HIGH & .643 & [.62,.67] \\
de & HIGH & .628 & [.60,.65] \\
ja & MID & .601 & [.57,.63] \\
hi & MID & .568 & [.54,.60] \\
bn & LOW & .518 & [.49,.55] \\
sw & LOW & .495 & [.46,.53] \\
\bottomrule
\end{tabular}
\caption{Llama~4 Scout zero-shot transfer \auroc{} (selected).}
\label{tab:scout-zs}
\end{table}

\subsection{CaCLR Is Inert}
\label{sec:caclr-results}

\begin{table}[t]
\centering
\small
\begin{tabular}{lcccl}
\toprule
\textbf{Method} & \textbf{HIGH} & \textbf{MID} & \textbf{LOW} & \textbf{Wins} \\
\midrule
Mean-shift & .000 & .000 & .000 & 0/19 \\
\caclr{} (best $\beta$) & .005 & .003 & $-$.008 & 8/19 \\
Lin.\ proj (50) & .025 & .015 & $-$.005 & 14/19 \\
Few-shot (25) & .038 & .030 & .022 & 17/19 \\
Few-shot (100) & .072 & .058 & .045 & 19/19 \\
\bottomrule
\end{tabular}
\caption{$\Delta$\auroc{} vs.\ zero-shot by transfer method and resource tier (Qwen3-235B, EN-optimal layer). \caclr{} provides negligible improvement; few-shot is the only consistently effective strategy.}
\label{tab:transfer-comparison}
\end{table}

Table~\ref{tab:transfer-comparison} presents the central negative result.
\caclr{}'s best configuration achieves $\Delta$\auroc{}$=0.002$ on average---within noise---while \emph{hurting} low-resource languages ($-$0.008).
Mean-shift produces exactly $\Delta$\auroc{}$=0.000$, as predicted by the rank-invariance proof under standardized logistic regression.

\paragraph{$\beta$ sweep.}
Table~\ref{tab:beta-sweep} shows that across three orders of magnitude ($\beta \in [0.001, 1.0]$), $\Delta$\auroc{} never exceeds noise level; all paired permutation tests yield $p{>}0.4$.
At $\beta{=}1.0$, the MMD penalty actively hurts ($\Delta{=}{-}0.003$).

\begin{table}[t]
\centering
\small
\begin{tabular}{ccccc}
\toprule
$\beta$ & $\Delta$\auroc{} & std & $\Delta$\ece{} & $p$ \\
\midrule
0.001 & 0.001 & 0.006 & 0.003 & 0.72 \\
0.01 & 0.002 & 0.007 & 0.005 & 0.58 \\
0.1 & 0.001 & 0.008 & 0.008 & 0.63 \\
1.0 & $-$0.003 & 0.011 & 0.012 & 0.41 \\
0 (BCE) & 0.000 & 0.005 & 0.001 & --- \\
\bottomrule
\end{tabular}
\caption{\caclr{} $\beta$ sweep (Qwen3-235B, avg.\ 20 languages). No $\beta$ achieves significant improvement over zero-shot.}
\label{tab:beta-sweep}
\end{table}

\paragraph{Component ablation.}
Table~\ref{tab:ablation} confirms no synergy between loss components.
BCE+SB-ECE provides marginal ECE improvement ($-$0.002); BCE+MMD provides marginal \auroc{} improvement (+0.002); the full combination equals the maximum of individual components, not their sum.

\begin{table}[t]
\centering
\small
\begin{tabular}{lccc}
\toprule
\textbf{Components} & $\Delta$\auroc{} & $\Delta$\ece{} & \textbf{Note} \\
\midrule
BCE only & .000 & .001 & baseline \\
BCE + SB-ECE & .001 & $-$.002 & $\alpha{=}0.1$ \\
BCE + MMD & .002 & .004 & $\beta{=}0.01$ \\
BCE + SB-ECE + MMD & .002 & .006 & no synergy \\
\bottomrule
\end{tabular}
\caption{\caclr{} component ablation (Qwen3-235B, avg.\ 20 languages).}
\label{tab:ablation}
\end{table}

\paragraph{Why does MMD fail?}
The MMD consistency loss operates on 1-D probe logits.
It matches marginal distributions $P(\hat{y}_{\text{en}})$ and $P(\hat{y}_t)$, but the core transfer problem is that conditional distributions $P(\text{correct} | \mathbf{h})$ differ across languages \cite{zhao2019learning}.
In MoE architectures, different languages activate different expert subsets, making the representation-space topology language-dependent in ways that 1-D marginal matching cannot address.
Linear projection succeeds (14/19 languages over \caclr{}) because it directly learns the subspace rotation rather than indirectly matching distributions.

\subsection{Capacity Artifacts in Cross-Lingual Probing}
\label{sec:control-results}

\begin{table}[t]
\centering
\small
\begin{tabular}{lcccc}
\toprule
\textbf{Tier} & \textbf{Shallow} & \textbf{Mid} & \textbf{Deep} & \textbf{Art.\ rate} \\
\midrule
EN (source) & 0.03 & 0.12 & 0.25 & 25\% \\
HIGH (4) & 0.03 & 0.09 & 0.18 & 32\% \\
MID (12) & 0.02 & 0.06 & 0.10 & 48\% \\
LOW (3) & 0.01 & 0.03 & 0.03 & 72\% \\
\bottomrule
\end{tabular}
\caption{Selectivity (= $\auroc{}_\text{real} - \auroc{}_\text{control}$) by resource tier and depth bin (Qwen3-235B). Artifact rate: fraction of lang--layer combinations with selectivity${<}0.05$.}
\label{tab:selectivity}
\end{table}

Table~\ref{tab:selectivity} presents the most consequential finding for the probing methodology.
Control probes trained on permuted labels achieve \auroc{} nearly as high as real-label probes for 72\% of low-resource language--layer combinations.
This means that for Swahili, Bengali, and Telugu, the entire probing pipeline produces results that are likely capacity artifacts rather than genuine correctness signal detection.

Table~\ref{tab:artifact-detail} provides per-language detail at the best layer.
Low-resource languages show selectivity of 0.02--0.03: real \auroc{} of 0.508--0.538 against control \auroc{} of 0.478--0.518.
This finding retroactively undermines Pareto frontier analysis for low-resource languages---their data points are artifact-suspect.

\begin{table}[t]
\centering
\small
\begin{tabular}{llcccc}
\toprule
\textbf{Lang} & \textbf{Tier} & \textbf{Sel.} & \textbf{Art?} & \textbf{Real} & \textbf{Ctrl} \\
\midrule
en & src & .25 & \texttimes & .760 & .510 \\
zh & HIGH & .18 & \texttimes & .672 & .492 \\
es & HIGH & .17 & \texttimes & .668 & .498 \\
ko & MID & .09 & \texttimes & .602 & .512 \\
hi & MID & .10 & \texttimes & .588 & .488 \\
tr & MID & .06 & $\sim$ & .571 & .511 \\
ar & MID & .05 & $\sim$ & .555 & .505 \\
sw & LOW & .03 & \checkmark & .508 & .478 \\
bn & LOW & .02 & \checkmark & .538 & .518 \\
te & LOW & .02 & \checkmark & .521 & .501 \\
\bottomrule
\end{tabular}
\caption{Per-language selectivity at the best layer (Qwen3-235B). Low-resource languages are flagged as artifacts; their real \auroc{} barely exceeds control \auroc{}.}
\label{tab:artifact-detail}
\end{table}

\subsection{CKA Partial Correlation Is Non-Significant}
\label{sec:cka-results}

\begin{table}[t]
\centering
\small
\begin{tabular}{lcccl}
\toprule
\textbf{Model} & \textbf{Part.\ $r$} & \textbf{$p$} & \textbf{Pool $\rho$} & \textbf{Interp.} \\
\midrule
Qwen3 & 0.04 & 0.58 & $-$0.28$^{**}$ & depth conf. \\
Scout & 0.08 & 0.35 & $-$0.35$^{**}$ & depth conf. \\
\bottomrule
\end{tabular}
\caption{CKA partial correlation controlling for depth. Pooled correlations are significant, but this is entirely driven by the depth confound: deeper layers simultaneously have lower CKA and worse transfer.}
\label{tab:cka-partial}
\end{table}

Table~\ref{tab:cka-partial} shows that naively correlating CKA with transfer degradation produces significant pooled correlations ($\rho{=}{-}0.28$, $p{=}0.008$ for Qwen3; $\rho{=}{-}0.35$, $p{=}0.003$ for Scout).
However, partial correlation controlling for normalized depth reduces these to non-significance ($r{=}0.04$, $p{=}0.58$; $r{=}0.08$, $p{=}0.35$).
This is a form of Simpson's paradox: depth simultaneously drives CKA decline and transfer degradation, creating a spurious direct correlation.

\paragraph{Same-layer analysis.}
At individual layers (holding depth constant), languages with higher CKA do show less transfer degradation ($\rho_{\text{avg}}{=}{-}0.42$), but only 5/36 layer comparisons across both models reach significance.
CKA is at best a weak between-language predictor at fixed depth, not the cross-depth explanatory variable initially hypothesized.

\paragraph{Subspace angles.}
The angle between English and target-language correctness directions increases with both depth and typological distance (Table~\ref{tab:subspace}).
Low-resource languages already show large angles (35--40\textdegree) at shallow layers, suggesting that their representations never meaningfully align with English on the correctness dimension.

\begin{table}[t]
\centering
\small
\begin{tabular}{llccc}
\toprule
\textbf{Lang} & \textbf{Tier} & \textbf{L48} & \textbf{L72} & \textbf{L88} \\
\midrule
zh & HIGH & 15\textdegree & 22\textdegree & 35\textdegree \\
es & HIGH & 14\textdegree & 20\textdegree & 32\textdegree \\
de & HIGH & 18\textdegree & 25\textdegree & 38\textdegree \\
ko & MID & 25\textdegree & 35\textdegree & 48\textdegree \\
hi & MID & 28\textdegree & 38\textdegree & 52\textdegree \\
sw & LOW & 38\textdegree & 48\textdegree & 58\textdegree \\
bn & LOW & 35\textdegree & 45\textdegree & 55\textdegree \\
te & LOW & 40\textdegree & 50\textdegree & 58\textdegree \\
\bottomrule
\end{tabular}
\caption{Subspace angle (degrees) between English and target-language correctness directions at representative layers (Qwen3-235B).}
\label{tab:subspace}
\end{table}

\subsection{Difficulty Decomposition Is Noisy}
\label{sec:decomp-results}

\begin{table}[t]
\centering
\small
\begin{tabular}{llcccc}
\toprule
\textbf{Lang} & \textbf{Tier} & \textbf{Diff.\ $\Delta$} & \textbf{Rand.\ $\Delta$} & \textbf{Gap} & $p$ \\
\midrule
zh & HIGH & .015 & $-$.012 & .027 & .031 \\
es & HIGH & .013 & $-$.014 & .027 & .028 \\
fr & HIGH & .011 & $-$.016 & .027 & .035 \\
de & HIGH & .008 & $-$.013 & .021 & .062 \\
hi & MID & .005 & $-$.015 & .020 & .088 \\
ru & MID & .003 & $-$.018 & .021 & .075 \\
ko & MID & .002 & $-$.014 & .016 & .145 \\
tr & MID & .001 & $-$.017 & .018 & .118 \\
sw & LOW & $-$.005 & $-$.019 & .014 & .212 \\
bn & LOW & $-$.010 & $-$.016 & .006 & .481 \\
\bottomrule
\end{tabular}
\caption{Difficulty decomposition: $\Delta$\auroc{} vs.\ no-projection baseline. Difficulty projection is better than random projection, but the effect is small and significant only for high-resource languages.}
\label{tab:decomp}
\end{table}

Table~\ref{tab:decomp} shows that projecting out the difficulty direction provides only a marginal average improvement of $\Delta$\auroc{}$=+0.005$---far smaller than few-shot (25) at $+0.032$.
The difficulty-vs-random gap is significant for only 4/10 languages, all high-resource.
IRT difficulty scores themselves are noisy ($R^2=0.35$ on average), and the ``difficulty-shared'' fraction estimate is highly unstable across languages (std$=0.18$), precluding a consistent decomposition narrative.
The difficulty component exists but is not cleanly separable in MoE architectures.

\subsection{Cross-Lingual Orthogonality Degrades for Low-Resource Languages}
\label{sec:ortho-results}

\begin{table}[t]
\centering
\small
\begin{tabular}{llccc}
\toprule
\textbf{Lang} & \textbf{Tier} & \textbf{Qwen3} & \textbf{Scout} & \textbf{Avg} \\
\midrule
en & src & 0.03 & 0.03 & 0.03 \\
zh & HIGH & 0.04 & 0.03 & 0.035 \\
es & HIGH & 0.03 & 0.04 & 0.035 \\
de & HIGH & 0.04 & 0.04 & 0.04 \\
ko & MID & 0.06 & 0.07 & 0.065 \\
hi & MID & 0.07 & 0.08 & 0.075 \\
tr & MID & 0.08 & 0.09 & 0.085 \\
bn & LOW & 0.12 & 0.14 & 0.13 \\
te & LOW & 0.15 & 0.13 & 0.14 \\
sw & LOW & 0.14 & 0.18 & 0.16 \\
\bottomrule
\end{tabular}
\caption{Cosine similarity between correctness and verbalized confidence directions. Orthogonality holds for high-resource languages (cosine${<}0.05$) but degrades for low-resource languages (up to 0.18), qualifying the ``universal orthogonality'' claim.}
\label{tab:ortho}
\end{table}

Table~\ref{tab:ortho} shows that the correctness--confidence orthogonality reported by Miao \& Ungar \citeyear{miao2026orthogonality} (cosine${<}0.04$ in English) holds for high-resource languages but systematically degrades for low-resource languages.
15\% of language--architecture combinations exceed cosine$=0.15$; all are low-resource languages.
This could reflect genuine structural non-orthogonality or, more likely given the control task results (\S\ref{sec:control-results}), noisy direction estimates from capacity-artifact probes.

\subsection{Linear Probes Transfer Better Than MLPs}
\label{sec:complexity}

\begin{table}[t]
\centering
\small
\begin{tabular}{lcccc}
\toprule
\textbf{Metric} & \textbf{LR} & \textbf{MLP} & $\Delta$ & $p$ \\
\midrule
EN \auroc{} & .760 & .785 & +.025 & ${<}.001$ \\
ZS avg (20) & .590 & .575 & $-$.015 & .042 \\
Degradation avg & .155 & .195 & +.040 & .008 \\
ECE avg (20) & .180 & .220 & +.040 & .015 \\
\bottomrule
\end{tabular}
\caption{LR vs.\ MLP probes (Qwen3-235B). MLP's English advantage is offset and reversed in cross-lingual transfer.}
\label{tab:lr-vs-mlp}
\end{table}

Table~\ref{tab:lr-vs-mlp} confirms that linear probes are the correct choice for cross-lingual transfer.
The MLP's 2.5pp English \auroc{} advantage is fully offset by 4pp excess transfer degradation, resulting in worse cross-lingual performance ($p{=}0.042$).
The MLP captures nonlinear, language-specific patterns (including expert-routing-dependent features in MoE architectures) that do not generalize.
LR transfers significantly better in 14/19 languages (bootstrap $p{<}0.05$).

\subsection{Linear Projection Sample Efficiency}
\label{sec:proj-results}

\begin{table}[t]
\centering
\small
\begin{tabular}{ccccc}
\toprule
$N$ & $\Delta$\auroc{} & std & Wins & Sig. \\
\midrule
10 & .005 & .032 & 10/19 & 2/19 \\
20 & .010 & .025 & 12/19 & 4/19 \\
50 & .015 & .020 & 14/19 & 7/19 \\
100 & .022 & .014 & 16/19 & 10/19 \\
200 & .028 & .010 & 17/19 & 13/19 \\
\bottomrule
\end{tabular}
\caption{Linear projection $\Delta$\auroc{} by number of parallel pairs (Qwen3-235B, avg.\ 20 languages).}
\label{tab:proj-efficiency}
\end{table}

\begin{table}[t]
\centering
\small
\begin{tabular}{lccc}
\toprule
\textbf{Tier} & $\Delta$\auroc{} & std & Sig. \\
\midrule
HIGH (4) & .025 & .012 & 3/4 \\
MID (13) & .015 & .018 & 4/13 \\
LOW (3) & $-$.005 & .028 & 0/3 \\
\bottomrule
\end{tabular}
\caption{Linear projection (50 pairs) by resource tier. Unstable and \emph{negative} for low-resource languages.}
\label{tab:proj-tier}
\end{table}

Tables~\ref{tab:proj-efficiency}--\ref{tab:proj-tier} show that linear projection with 50 parallel pairs outperforms \caclr{} on 14/19 languages but fails for low-resource languages.
With 50 pairs, low-resource languages show negative average $\Delta$ ($-$0.005) with high variance (std$=$0.028); even at 200 pairs, Swahili, Bengali, and Telugu do not consistently improve.
This corroborates the subspace angle analysis: low-resource correctness directions diverge too far from English for low-rank linear alignment to help.

\subsection{Cross-Benchmark Generalization}
\label{sec:generalization}

\begin{table}[t]
\centering
\small
\begin{tabular}{lcccc}
\toprule
\textbf{Bench.} & \textbf{EN} & \textbf{ZS} & \textbf{CaCLR $\Delta$} & \textbf{Rank $r$} \\
\midrule
GlobalMMLU & .76 & .59 & .002 & --- \\
Belebele & .72 & .56 & .001 & .82 \\
MMLU-ProX & .74 & .57 & .002 & .88 \\
\bottomrule
\end{tabular}
\caption{Cross-benchmark generalization (Qwen3-235B). Language ranking is highly consistent (Kendall $\tau{=}0.78$); \caclr{} is consistently inert.}
\label{tab:generalization}
\end{table}

Table~\ref{tab:generalization} shows that findings generalize across benchmarks.
Language rankings are highly consistent with GlobalMMLU (Belebele rank $r{=}0.82$; MMLU-ProX $r{=}0.88$; Kendall $\tau{=}0.78$), and \caclr{} remains inert on both alternative benchmarks.
The low-resource \auroc{} drop is consistent across all three benchmarks, confirming that the viability boundary is task-general, not a GlobalMMLU-specific artifact.

\subsection{Quantization Validation}
\label{sec:quant}

\begin{table}[t]
\centering
\small
\begin{tabular}{llcccc}
\toprule
\textbf{Model} & \textbf{Q} & \textbf{Layer} & \textbf{FP16} & \textbf{Q} & $\Delta$ \\
\midrule
Scout & W4A16 & L8 & .645 & .642 & $-$.003 \\
Scout & W4A16 & L40 & .740 & .725 & $-$.015 \\
Qwen3 & INT4 & L16 & .640 & .635 & $-$.005 \\
Qwen3 & INT4 & L72 & .762 & .742 & $-$.020 \\
Qwen3 & INT4 & L88 & .745 & .723 & $-$.022 \\
\bottomrule
\end{tabular}
\caption{Quantization impact on probe \auroc{}. Deeper layers suffer larger degradation, but rank agreement exceeds 96\%.}
\label{tab:quant}
\end{table}

Table~\ref{tab:quant} confirms that quantization introduces a small systematic \auroc{} reduction (avg.\ $-$0.008 for Scout W4A16; $-$0.012 for Qwen3 INT4) that is larger at deeper layers but does not alter per-item rankings (rank agreement $>$96\%).
The $\sim$1--2pp quantization artifact is far smaller than the 10--25pp cross-lingual transfer degradation, and does not affect any qualitative conclusion.

\subsection{Extended Language Coverage}
\label{sec:extended}

\begin{table}[t]
\centering
\small
\begin{tabular}{lccc}
\toprule
\textbf{Category} & $n$ & \textbf{Criterion} & \textbf{Examples} \\
\midrule
Viable & 18 & \auroc{}${>}.60$ + sel${>}.05$ & en, zh, es, fr, ja, \ldots \\
Marginal & 12 & \auroc{}${>}.55$ + sel .03--.05 & hi, tr, pl, ar, \ldots \\
Non-viable & 22 & \auroc{}${<}.55$ or sel${<}.03$ & sw, bn, te, am, \ldots \\
\bottomrule
\end{tabular}
\caption{Viability classification across 52 languages. Only 18/51 target languages (35\%) meet joint viability criteria for probe-based calibration.}
\label{tab:viability}
\end{table}

Extending to 52 languages (Table~\ref{tab:viability}), we classify each target language by joint viability criteria: \auroc{}${>}0.60$ \emph{and} selectivity${>}0.05$.
Only 18/51 (35\%) qualify as viable---all high- or mid-resource.
22 languages (43\%) are non-viable, with artifact rates exceeding 85\% for newly added low-resource languages.
The viability boundary is sharp: no low-resource language passes both criteria.

%------------------------------------------------------------
\section{Discussion}
\label{sec:discussion}

\paragraph{Viability boundary interpretation.}
The 35\% viability rate should not be interpreted as an indictment of the probe-based calibration approach per se, but rather as a boundary condition that the field must acknowledge.
For the 18 viable languages, cross-lingual probe transfer is a legitimate calibration strategy with practical value.
For the remaining languages, alternative approaches are needed.

\paragraph{MoE vs.\ dense architectures.}
The collapse of the Pareto frontier on MoE architectures is mechanistically interpretable: input-dependent expert routing creates language-specific computational pathways that break the assumption of shared representational geometry underlying cross-lingual transfer.
As MoE becomes the dominant architecture for frontier models, this finding has direct implications for the scalability of probe-based methods.

\paragraph{Capacity artifacts as a field-level concern.}
The 72\% artifact rate for low-resource languages extends Hewitt \& Liang's \citeyear{hewitt2019designing} insight to a setting they did not consider.
Their original work focused on English syntactic probing; we show that capacity artifacts are dramatically amplified in multilingual settings, particularly for languages with lower-quality evaluation data (GlobalMMLU translation artifacts may contribute systematic patterns that probes exploit).
Any future cross-lingual probing study should include control task analysis as a mandatory validation step.

\paragraph{CKA methodology warning.}
The depth confound in CKA analysis (significant pooled correlation but non-significant partial correlation) is a cautionary example for the representation similarity literature.
Studies correlating CKA with downstream metrics without controlling for depth may report spurious relationships driven by the shared depth confounder.

\paragraph{Constructive paths forward.}
Three directions emerge for languages outside the viability boundary:
(1)~\textbf{Few-shot target-language probes} (25--100 examples) consistently improve transfer even for low-resource languages ($\Delta$\auroc{}$=+0.022$ at $k{=}25$ for low-resource, $+0.045$ at $k{=}100$), representing the most practical near-term solution.
(2)~\textbf{Multilingual probe training} on pooled high-resource languages may improve the shared correctness subspace without requiring target-language labels.
(3)~\textbf{Dense model selection}: when Pareto analysis is required, dense architectures should be preferred over MoE for cross-lingual probing applications.

\paragraph{Limitations.}
Our evaluation is restricted to MCQA tasks, which enable clean correctness labeling but may not represent all calibration scenarios.
Quantization introduces 1--2pp \auroc{} degradation that we account for but cannot eliminate.
The causal attribution of MoE routing to Pareto collapse is correlational (n$=$2 MoE vs.\ n$=$1 dense); future work should test on additional architectures.
Low-resource non-viability may partly reflect GlobalMMLU translation quality rather than representational limitations.

%------------------------------------------------------------
\section{Conclusion}
\label{sec:conclusion}

We presented the first large-scale evaluation of cross-lingual correctness probe transfer for LLM calibration, spanning 52 languages and three architectures.
Our findings establish a viability boundary: probe-based calibration transfers reliably for only 35\% of tested languages, all high- or mid-resource.
The depth--transferability Pareto frontier is architecture-dependent, collapsing on MoE models.
MMD-based distributional matching (\caclr{}) is inert across all configurations.
Most critically, control task analysis reveals that 38\% of cross-lingual probing results---and 72\% for low-resource languages---are capacity artifacts, extending Hewitt \& Liang's framework to multilingual settings as a field-level methodological warning.
Few-shot target-language probes (25--100 examples) remain the most reliable path forward for languages outside the viability boundary.

\begin{thebibliography}{30}
\bibitem[\protect\citeauthoryear{Ahuja et~al.}{2024}]{ahuja2024megaverse}
Ahuja, K., et~al. (2024).
\newblock Megaverse: Benchmarking large language models across languages, modalities, models and tasks.
\newblock In \emph{NeurIPS Datasets and Benchmarks}.

\bibitem[\protect\citeauthoryear{Angelopoulos and Bates}{2023}]{angelopoulos2023conformal}
Angelopoulos, A.~N. and Bates, S. (2023).
\newblock Conformal prediction: A gentle introduction.
\newblock \emph{Foundations and Trends in Machine Learning}.

\bibitem[\protect\citeauthoryear{Artetxe et~al.}{2018}]{artetxe2018robust}
Artetxe, M., Labaka, G., and Agirre, E. (2018).
\newblock A robust self-learning method for fully unsupervised cross-lingual mappings of word embeddings.
\newblock In \emph{ACL}.

\bibitem[\protect\citeauthoryear{Ben-David et~al.}{2010}]{bendavid2010theory}
Ben-David, S., Blitzer, J., Crammer, K., Kulesza, A., Pereira, F., and Vaughan, J.~W. (2010).
\newblock A theory of learning from different domains.
\newblock \emph{Machine Learning}, 79(1-2):151--175.

\bibitem[\protect\citeauthoryear{Burns et~al.}{2023}]{burns2023discovering}
Burns, C., Ye, H., Klein, D., and Steinhardt, J. (2023).
\newblock Discovering latent knowledge in language models without supervision.
\newblock In \emph{ICLR}.

\bibitem[\protect\citeauthoryear{Conneau et~al.}{2018}]{conneau2018word}
Conneau, A., Lample, G., Ranzato, M., Denoyer, L., and J{\'e}gou, H. (2018).
\newblock Word translation without parallel data.
\newblock In \emph{ICLR}.

\bibitem[\protect\citeauthoryear{Hewitt and Liang}{2019}]{hewitt2019designing}
Hewitt, J. and Liang, P. (2019).
\newblock Designing and interpreting probes with control tasks.
\newblock In \emph{EMNLP}.

\bibitem[\protect\citeauthoryear{Kornblith et~al.}{2019}]{kornblith2019similarity}
Kornblith, S., Norouzi, M., Lee, H., and Hinton, G. (2019).
\newblock Similarity of neural network representations revisited.
\newblock In \emph{ICML}.

\bibitem[\protect\citeauthoryear{Miao et~al.}{2026}]{miao2026coral}
Miao, N., et~al. (2026).
\newblock CORAL: Calibration via representation-aligned learning.
\newblock In \emph{NAACL}.

\bibitem[\protect\citeauthoryear{Miao and Ungar}{2026}]{miao2026orthogonality}
Miao, N. and Ungar, L. (2026).
\newblock Orthogonality of correctness and confidence in language model representations.
\newblock \emph{Preprint}.

\bibitem[\protect\citeauthoryear{Pires et~al.}{2019}]{pires2019multilingual}
Pires, T., Schlinger, E., and Garrette, D. (2019).
\newblock How multilingual is multilingual BERT?
\newblock In \emph{ACL}.

\bibitem[\protect\citeauthoryear{Voita and Titov}{2020}]{voita2020information}
Voita, E. and Titov, I. (2020).
\newblock Information-theoretic probing with minimum description length.
\newblock In \emph{EMNLP}.

\bibitem[\protect\citeauthoryear{Wu and Dredze}{2019}]{wu2019beto}
Wu, S. and Dredze, M. (2019).
\newblock Beto, Bentz, Becas: The surprising cross-lingual effectiveness of BERT.
\newblock In \emph{EMNLP}.

\bibitem[\protect\citeauthoryear{Xiao et~al.}{2025}]{xiao2025knowledge}
Xiao, Y., et~al. (2025).
\newblock Cross-lingual knowledge boundary probe transfer.
\newblock In \emph{ACL}.

\bibitem[\protect\citeauthoryear{Yu et~al.}{2025}]{yu2025confidence}
Yu, Z., et~al. (2025).
\newblock The confidence manifold in LLM representations.
\newblock In \emph{ICML}.

\bibitem[\protect\citeauthoryear{Zhao et~al.}{2019}]{zhao2019learning}
Zhao, H., Zhang, S., Wu, G., Moura, J.~M., Costeira, J.~P., and Gordon, G.~J. (2019).
\newblock On learning invariant representations for domain adaptation.
\newblock In \emph{ICML}.
\end{thebibliography}

\appendix

\section{Full Zero-Shot Transfer Results}
\label{app:full-zs}

Due to space constraints, Tables~\ref{tab:qwen-zs} and \ref{tab:scout-zs} show selected languages.
Full results for all 20 core languages on both architectures are available in the supplementary material.
On Qwen3-235B, the zero-shot transfer success rate (defined as \auroc{}${>}0.60$) is 55\% (11/20 languages).
On Llama~4 Scout, the rate drops to 40\% (8/20 languages).

\section{Dense Architecture (Gemma~4 31B) Details}
\label{app:gemma}

Gemma~4 31B Dense achieves EN \auroc{}$=0.78$ (L48, 80\% depth) with ZS transfer average of 0.61 across 20 languages.
Selectivity artifact rate is 22\%, substantially lower than both MoE architectures (38--41\%).
Linear projection with 50 pairs achieves $\Delta$\auroc{}$=0.028$ on average, nearly double the MoE rate.
These results confirm that dense architectures remain the favorable setting for cross-lingual probe transfer.

\end{document}
